% Part: proof-theory
% Chapter: normalization
% Section: introduction

\documentclass[../../../include/open-logic-section]{subfiles}

\begin{document}

\olfileid{pt}{nor}{seg}

\olsection{খণ্ড ও কর্তন}

!!^a{introduction} বিধির পরে !!a{elimination} বিধি এলে
তা !!a{proof}-এর একটি ঘুরপথ;
!!a{proof} স্বাভাবিকীকরণ মানে এমন ঘুরপথগুলি দূর করা।
কিন্তু \Elim{\lor} ও \Elim{\lexists}-এর জন্য
``যে ঘুরপথ দূর করতে চাই'' তার সংজ্ঞাই কিছুটা জটিল।
কারণ, ঘুরপথে !!{elimination} বিধি যে
!!{introduction} বিধির \emph{ঠিক পরেই} থাকবে,
তা নিশ্চিত নয়। মাঝখানে \Elim{\lor} বা
\Elim{\lexists} থাকতে পারে, যেমন:
\[
\AxiomC{}
\DeduceC{$\lexists[x][!C(x)]$}
\AxiomC{$\Discharge{!A}{x}$}
\DeduceC{$!B$}
\DischargeRule{\Intro{\lif}}{x}
\UnaryInfC{$!A \lif !B$}
\DischargeRule{\Elim{\lexists}}{y}
\BinaryInfC{$!A \lif !B$}
\AxiomC{}
\DeduceC{$!A$}
\RightLabel{\Elim{\lif}}
\BinaryInfC{$!B$}
\DisplayProof
\]
বস্তুত, যেকোনো সংখ্যক \Elim{\lor} বা
\Elim{\lexists} \Intro{\lif} ও \Elim{\lif}-কে
পৃথক করতে পারে, যেমন:
\[
\AxiomC{}
\DeduceC{$\lexists[x][!C(x)]$}
\AxiomC{}
\DeduceC{$!D \lor !E$}
\AxiomC{$\Discharge{!A}{x}$}
\DeduceC{$!B$}
\DischargeRule{\Intro{\lif}}{x}
\UnaryInfC{$!A \lif !B$}
\AxiomC{$\Discharge{!A}{x}$}
\DeduceC{$!B$}
\DischargeRule{\Intro{\lif}}{x}
\UnaryInfC{$!A \lif !B$}
\RightLabel{\Elim{\lor}}
\TrinaryInfC{$!A \lif !B$}
\DischargeRule{\Elim{\lexists}}{y}
\BinaryInfC{$!A \lif !B$}
\AxiomC{}
\DeduceC{$!A$}
\RightLabel{\Elim{\lif}}
\BinaryInfC{$!B$}
\DisplayProof
\]
এই উদাহরণে একই \Elim{\lif} একাধিক ঘুরপথে
থাকতে পারে: এর আগে একাধিক~\Intro{\lif} আছে।
সব সম্ভাবনা ধরতে \Log{N1i}-প্রমাণে
!!{formula}-সংঘটনের বিশেষ ধরনের অনুক্রম
দিয়ে ঘুরপথ সংজ্ঞায়িত করি; এগুলিকে
\emph{খণ্ড} বলি। এখানে $!A \lif !B$-এর
এমন সংঘটনগুলির অনুক্রম নেওয়া হয়,
যেখানে $\Elim{\lor}$ বা $\Elim{\lexists}$-এর
কোনো গৌণ পূর্বধারণার পরে ওই বিধির সিদ্ধান্ত আসে।
উদাহরণে এমন দুটি খণ্ড আছে; প্রতিটিতে
$!A \lif !B$-এর তিনটি সংঘটন।
আনুষ্ঠানিক সংজ্ঞাটি এই:

\begin{defn}
\Log{N1i} !!{proof}~$\delta$-তে
\emph{খণ্ড} হলো ওই~$\delta$-তে একই !!{formula}~$!A$-এর
সংঘটন $!A_1$, \dots, $!A_n$-এর
এমন অনুক্রম, যেখানে
\begin{enumerate}
  \item $!A_1$ \Elim{\lor} বা \Elim{\lexists}-এর
  সিদ্ধান্ত নয়;
  \item $!A_n$ \Elim{\lor} বা \Elim{\lexists}-এর
  গৌণ পূর্বধারণা নয়;
  \item $i = 1$, \dots,~$n-1$ হলে,
  $!A_i$ \Elim{\lor} বা \Elim{\lexists}-এর
  একটি গৌণ পূর্বধারণা এবং $!A_{i+1}$
  ওই প্রয়োগের সিদ্ধান্ত।
\end{enumerate}
খণ্ডটির \emph{দৈর্ঘ্য} হলো~$n$।
\end{defn}
% BN-SRC-582 TARGET-CORRECTION-BEGIN
\par\smallskip\noindent\textnormal{\small\textbf{উৎস-সংশোধন (BN-SRC-582):} খণ্ডের সংজ্ঞায় একই $!A$-সূত্রের পরবর্তী সংঘটন উৎসের $A_{i+1}$ নয়, $!A_{i+1}$; বাদ পড়া সূত্রচিহ্ন ফিরেছে।}
% BN-SRC-582 TARGET-CORRECTION-END

সব খণ্ডই দূর করার মতো ঘুরপথ নয়।
যেগুলি ঘুরপথ, সেগুলিকে \emph{কর্তন} বলি।

\begin{defn}
\emph{কর্তন} হলো এমন খণ্ড, যাতে $!A_n$
কোনো !!a{elimination} প্রয়োগের প্রধান পূর্বধারণা
এবং হয় $n=1$ ও $!A_1 = !A_n$ কোনো
!!a{introduction} প্রয়োগের বা~\FalseInt-এর
সিদ্ধান্ত, নয়তো~$n>1$।
\end{defn}

লক্ষ করুন, কর্তন-খণ্ডের শুরুতে
!!a{introduction} বিধির সিদ্ধান্ত থাকতেই হবে না।
খণ্ডটির শেষে !!a{elimination} বিধির
প্রধান পূর্বধারণা থাকলেই এই শর্ত পূরণ হয়।
তবে দৈর্ঘ্য~$1$-এর খণ্ড কর্তন হতে হলে
সেটি !!{introduction} বিধি অথবা
মিথ্যা-ভূমিকা বিধির সিদ্ধান্তও হতে হবে।
% BN-SRC-583 TARGET-CORRECTION-BEGIN
\par\smallskip\noindent\textnormal{\small\textbf{উৎস-সংশোধন (BN-SRC-583):} ঠিক আগের কর্তন-সংজ্ঞায় দৈর্ঘ্য একের ক্ষেত্রে \FalseInt-এর সিদ্ধান্তও অনুমোদিত; উৎসের শেষ ব্যাখ্যাবাক্যে ওই বিকল্পটি বাদ ছিল।}
% BN-SRC-583 TARGET-CORRECTION-END

\begin{defn}
কর্তনের \emph{মাত্রা} হলো~$\depth{!A}$।
!!a{proof}~$\delta$-র \emph{কর্তনমাত্রা}
~$\cutrank{\delta}$ হলো $\delta$-তে
কর্তনগুলির মধ্যে সর্বোচ্চ মাত্রা।
কোনো কর্তনের মাত্রা $\cutrank{\delta}$
হলে সেটি $\delta$-তে \emph{সর্বাধিক}
কর্তন। $\delta$-র \emph{কর্তনদৈর্ঘ্য}
হলো তার সব সর্বাধিক কর্তনের দৈর্ঘ্যের যোগফল।
\end{defn}

\end{document}
